\subsection{Discussion}
\label{rasc:sec:discussion}

\rev{\noindent\textbf{Action-length predictability and richer observability.}} Action lengths of some devices (e.g., shades) are more predictable than others, e.g., a thermostat  (HVAC) may have seasonal variations, oven cook time depends on food, etc. \Cref{rasc:sec:eval} evaluated {\it representative} devices that span this spectrum. 
More expressive observability (e.g., progress bars) may be feasible atop \abs{}, but it must be {\it simple} and {\it accurate}.

\rev{\noindent\textbf{Causality expressiveness.}} Current automation interfaces (e.g., Alexa~\cite{Alexa}, Google Home~\cite{GoogleHome}) support limited causality expressiveness.
Their IFTTT-style~\cite{ifttt} rules struggle with conditional, multi-stage dependencies.
Future work could integrate expressive, flow-based (and visual) systems like Node-RED~\cite{NodeRED} or  Flogo~\cite{flogo}.

\rev{\noindent\textbf{Privacy.} \abs{} does not alter the privacy posture of the deployment it runs in. All of its polling, lifecycle-state tracking, and scheduling operate within the same administrative domain as the underlying platform (e.g., a home running Home Assistant~\cite{HomeAssistant}): \abs{} introduces no new data flows that cross the domain boundary, sends no telemetry to third parties, and observes only device state that the platform already collects. A deployment that keeps its control logic on-premise therefore retains exactly the same privacy boundary with \abs{} as without it.}

\rev{\noindent\textbf{Scaling to many users.} Our evaluation considers routines issued by a single household's automation platform; larger deployments (e.g., a commercial building) may have many users concurrently issuing actions and routines. Scaling \abs{} to this regime is a natural direction for future work. Because lifecycle states are per-device rather than per-user, concurrent requests that touch the same device can share a single polling stream: \emph{batching} status queries across users' actions, and serving all interested routines from one shared lifecycle-state cache, would keep polling overhead proportional to the number of active devices rather than the number of users.}

\rev{\noindent\textbf{Maintaining conditions over long periods.} Some automation goals are not one-shot but \emph{standing}---e.g., ``the floor should stay clean,'' where a pet may drop food repeatedly throughout the day. In today's edge-triggered platforms, every sensed instance fires the routine anew, even if a cleaning run is already underway. \abs{}'s lifecycle abstraction offers a principled remedy: because the platform can observe that an instance of the same action is currently between \textsc{Start} and \textsc{Complete}, it can suppress redundant firings while the condition is being restored, and the \textsc{Complete} event naturally re-arms the trigger for the next violation. Standing goals thus map onto \abs{} as recurring actions with observable progress, turning duplicate-trigger suppression from a per-routine hack into a property the abstraction provides.}

\rev{\noindent\textbf{Exposing tradeoffs and costs to users.} \abs{}'s user tolerance threshold $Q_w$ (\Cref{rasc:sec:rasc}) already encodes a tradeoff between detection latency and polling cost, but it is currently a single, fixed value. Surfacing this tradeoff to users---and letting them set context-dependent thresholds, e.g., different $Q_w$ values for night versus day, in whichever direction suits the household---would let occupants tune responsiveness against network and energy overhead, much as they tune thermostat schedules today.}
